% 原文 pp.173–180（印刷 pp.1291–1298）
\section{Lecture 13：二维 $\mathcal{N}=2$ 超对称理论（一）}

\begin{keytakeaways}
本讲研究二维 $\mathcal{N}=2$ 超对称场论，重点是带复值手征超场的 $U(1)$ 规范理论：这类理论由两个函数——超势 $W$ 与扭曲超势 $\widetilde W$——决定，前者是超场的复多项式（须满足「权」条件），后者由单个复参数 $t=-ir+\theta/2\pi$（更准确地说由 $\exp(-2\pi it)$）决定，故理论族由柱面 $\mathbb{C}/\mathbb{Z}$ 参数化。物理图像是柱面两端各有一个极限：$\operatorname{Im}t\to-\infty$ 一端是靶空间为射影簇 $X$ 的 $\sigma$ 模型，$\operatorname{Im}t\to+\infty$ 一端是 Landau--Ginzburg 模型的轨形版本，中间由重整化群流连接。数学工具是 Kähler 几何与指标定理：$\beta$ 函数被看成 $X$ 上的闭 $(1,1)$-形式，其零点方程 $R=\bar\partial\partial U$ 由 Yau 定理求解；反常则由 $c_1(X)$ 在 Riemann 面上的积分给出。主要结论有三：$c_1(X)=0$ 时 $\beta$ 函数对所有 $t$ 为零，得到一族共形理论；$c_1(X)\neq0$ 时存在非平凡的重整化群流（$c_1>0$ 使 $\operatorname{Im}t$ 增大，$c_1<0$ 使其减小）；要求 $R$ 对称性无反常则须 $\sum_iq_i=0$，这正是本讲最后所取的权 $q=(1,\dots,1,-n)$ 的由来。
\end{keytakeaways}

% ===== 原文 p.173 =====
\subsection{引言}
% 原文 §13.1：Introduction

在本讲以及接下来的两讲中，我们讨论二维的 $\mathcal{N}=2$ 超对称场论。\footnote{本讲笔记由 John Morgan 与 David Morrison 整理。（原文 p.1291 页脚）} 今天，在就这类理论作一些一般性的介绍性评述之后，我们考虑二维的 $U(1)$ 规范理论，其中的手征超场取复值、并在规范群下带荷。这类理论依赖于两个函数——一个超势（superpotential）$W$ 与一个扭曲超势（twisted superpotential）$\widetilde W$。超势 $W$ 是超场的一个复多项式。对于今天将要关心的一类模型，这个多项式必须满足某个「权」（weight）条件，除此之外则是任意的。

对规范群 $U(1)$ 而言，我们所考察的那一类模型的扭曲超势 $\widetilde W$ 由单个复参数
\[
t=-ir+\frac{\theta}{2\pi},
\]
或者说更确切地由 $\exp(-2\pi it)$ 决定。今天我们保持超势固定，而让扭曲超势变动。于是，对每个满足权条件的（一般的）复多项式 $W$，我们都得到一族理论，它们由柱面 $\mathbb{C}/\mathbb{Z}$ 上的点 $\exp(-2\pi it)$ 参数化。我们关心的是柱面两端 $\operatorname{Im}(t)\to+\infty$ 与 $\operatorname{Im}(t)\to-\infty$ 处会发生什么。在我们所考虑的情形中，$\operatorname{Im}(t)\to-\infty$ 的极限可以用一个 $\sigma$ 模型来描述，其靶空间是由给出超势的多项式 $W$ 所定义的射影簇（projective variety）$X$。在这个区域里，$-\operatorname{Im}(t)$ 就是 $X$ 的 Kähler 类；而超势中只有一个参数这件事，与 $H^{1,1}(X)$（至少是其中从某个环绕的加权射影空间（weighted projective space）拉回的那一部分）是一维的这件事相关联。$\operatorname{Im}(t)\to+\infty$ 的极限则用一个 Landau--Ginzburg 模型的轨形（orbifold）版本来描述。

当 $c_1(X)=0$ 时，我们将发现 $\beta$ 函数对所有的 $t$ 都为零，从而得到一族共形不变的理论。当 $c_1(X)\neq0$ 时，我们得到的则是 $t$ 中的非平凡重整化群流。当 $c_1(X)>0$ 时这个流使 $\operatorname{Im}(t)$ 增大，当 $c_1(X)<0$ 时则使其减小。（此外，在后两种情形中，柱面端点处还有「额外的」真空，红外下将流动到那里。）这使我们能把各种 Landau--Ginzburg 模型看作 $\sigma$ 模型的红外极限，并把另一些 $\sigma$ 模型看作 Landau--Ginzburg 模型的红外极限。这之所以有意思，是因为这两类模型在经典层面的描述非常不同，

% [Tutor Note: 这里的参数约定是全文的枢纽，务必按原文读作 $\operatorname{Im}t=-r$。于是 $\operatorname{Im}t\to-\infty$ 等价于 $r\to+\infty$，即 Fayet--Iliopoulos 参数（也就是 Kähler 类，见 p.175）非常大——这正是 $\sigma$ 模型描述成立的大体积极限；初学者常把「$t\to-\infty$」误当成弱耦合极限，从而把柱面的两端弄反。又，$\exp(-2\pi it)=e^{-2\pi r}e^{-i\theta}$ 才是良定义的坐标：$\theta$ 角只有 $2\pi$ 的周期性（$\theta\to\theta+2\pi$ 不改变理论），故 $t\to t+1$ 是等价关系，参数空间才是柱面 $\mathbb{C}/\mathbb{Z}$ 而非复平面。柱面上「被打洞的那一点」即 $\exp(-2\pi it)=0$，对应 $r\to+\infty$ 的极限点，也就是 $\sigma$ 模型那一端；这解释了 p.175 所说「理论在该点有某种极点」。]

% ===== 原文 p.174 =====
在经典层面即是如此。今天讲的大部分内容将处理一些已被广泛研究的特殊情形（$c_1(X)=0$ 的情形，以及某些 $c_1(X)>0$、流向有质量红外理论的情形），但本讲末尾将要给出的总体图景，据我所知实际上是新的。

\subsection{二维 $\mathcal{N}=2$ 超对称理论的一般性讨论}
% 原文 §13.2：Generalities on $\mathcal{N}=2$ SUSY theories in dimension two

\subsubsection{$\mathcal{N}=2$ 超对称 $\sigma$ 模型的 $\beta$ 函数}
% 原文 §13.2.1：The $\beta$-function of $\mathcal{N}=2$ SUSY sigma models

第一条评述涉及靶空间为紧 Kähler 流形 $X$ 的 $\mathcal{N}=2$ 超对称 $\sigma$ 模型，我们假定 $X$ 具有 Kähler--Einstein 度量。在 [II-CFT, §3.8]（式 3.42）中我们已经看到，单圈 $\beta$ 函数由度量的 Ricci 曲率的符号决定；事实上，当 Ricci 张量正定或负定时，单圈 $\beta$ 函数的符号与 Ricci 曲率的符号相反。因此，若 $c_1(X)>0$，则单圈 $\beta$ 函数为负。（所谓 $c_1(X)$，我们指 $X$ 的切丛的 $c_1$——所以对复射影空间有 $c_1>0$。）在这种情形下，小耦合时 $\beta$ 函数为负，因此所考虑的 $\sigma$ 模型是一个「好的」基本理论，它是渐近自由（asymptotically free）的。这样的理论在量子层面是良定义的。于是问题来了：它在红外流向何处？今天我们将看到这样一些 $X$ 的例子，其红外理论是有质量且平凡的；也会看到流向非平凡红外不动点的例子。在后一情形中，我们所得到的极限超共形理论可以用所谓的超对称 Landau--Ginzburg 理论来描述。Landau--Ginzburg 模型有时可以用代数方法显式求解。

另一方面，若 $c_1(X)<0$，则单圈 $\beta$ 函数为正，因而在弱耦合时 $\beta$ 为正。于是我们没有理由指望这样的 $X$ 上的 $\sigma$ 模型是「好的」基本理论。相反，在这种情形下该 $\sigma$ 模型是一个有效理论，在红外是自由的；人们应当问：能否描述一个具体的、良定义的理论，使它在红外流向这些 $\sigma$ 模型？我们将通过取 $\operatorname{Im}(t)\to\infty$ 来对这类 $\sigma$ 模型解决这个问题：在那里我们会找到一个 Landau--Ginzburg 理论（它是超共形的），它经过一个相关微扰（relevant perturbation，即把 $\operatorname{Im}(t)$ 取大而正、但不为 $\infty$）流向靶空间为 $X$ 的 $\sigma$ 模型（外加一些额外的有质量真空）。

$c_1(X)=0$ 的情形，即 $X$ 为 Calabi--Yau 的情形，则不同。此时单圈 $\beta$ 函数为零，单圈贡献不改变 Kähler 度量。但结果表明 $\beta$ 函数还有更高圈的修正，所以它并不恒为零，因而 $X$ 上（带 Ricci 平坦度量）的 $\sigma$ 模型并非共形不变的。其结构实际上如下。$X$ 上的 Kähler 度量自然与一个闭 $(1,1)$-形式 $\omega$（即 Kähler 形式）相联系。重整化群方程是关于 $\omega$ 相对于重整化群标度 $\mu$ 的变化的方程，其一般形式为
\[
\frac{d\omega}{d\ln\mu}=\beta,
\]
其中 $\beta$「函数」最好像 $\omega$ 一样理解成 $X$ 上的闭 $(1,1)$-形式。现在，$\beta$ 的 $s$ 圈贡献有如下特征：它是 $X$ 的 Kähler 几何的局部泛函；更具体地说，若我们定义一种分次（grading），把导数计为 $1$、把 Riemann 张量计为 $2$，则 $s$ 圈部分

% [Tutor Note: 「单圈 $\beta$ 函数的符号与 Ricci 曲率相反」是本讲的起点，也是最容易因约定不同而整体差一个符号的地方：它取决于世界面动能项的归一化（$\frac{1}{4\pi\alpha'}\int g_{i\bar j}\partial\phi^i\bar\partial\phi^{\bar j}$ 与相反符号的写法会给出相反的 $\beta$）。一个可以自查的物理判据是：$\mathbb{CP}^1$ 的 $\sigma$ 模型是众所周知的渐近自由理论，而 $\mathbb{CP}^1$ 的 Ricci 曲率为正，故「$c_1>0\Rightarrow\beta<0$」自洽；反之 $c_1<0$（例如亏格 $\ge2$ 的曲线作为靶空间）时弱耦合下 $\beta>0$，理论不是紫外完备的，只能当作红外有效理论——这正是作者随后要用 Landau--Ginzburg 模型为其补上紫外描述的原因。]

% [Tutor Note: 把 $\beta$ 当成 $X$ 上的闭 $(1,1)$-形式、并写出 $d\omega/d\ln\mu=\beta$，是本节的关键一步：Kähler 度量的形变在 $(1,1)$ 类中取值，而重整化群流只在 $\omega$ 的 Kähler 类与复结构这两个方向上起作用。注意这里 $\mu$ 取的是「标度」而非动量，故是 $d/d\ln\mu$ 而非 $\mu\,d/d\mu$；p.175 中作者又写成 $\mu\,d\omega/d\mu=0$，两式只差 $\omega$ 的一个整体因子，指同一个零点条件。]

% ===== 原文 p.175 =====
（接上页）的 $\beta$ 函数次数为 $2s$。\footnote{$\beta$ 是局部的，因为它由 Feynman 图算出；Feynman 图中的所有顶点在 $X$ 上都是局部的。关于 $\beta$ 的 $s$ 圈贡献之分次的陈述，是通过考察图中各顶点的权来证明的；越来越高阶的顶点来自 $X$ 的度量在正规坐标中展开的高阶项，因而具有更高的权。} 对 $s=1$，$X$ 的 Ricci 张量具有正确的性质，而单圈计算表明 $\beta$ 中的单圈项是 Ricci 张量 $R$ 的一个负的倍数。$s>1$ 时又如何呢？标准的不变量理论表明：Kähler 流形上任何可以局部地、普适地写出的、次数为 $2s$（$s>1$）的闭 $(1,1)$-形式，都具有 $\bar\partial\partial U$ 的形式，其中 $U$ 是局部构造出来的、次数为 $2s-2$ 的函数。因此，$\beta$ 函数零点所满足的方程（$\mu\,d\omega/d\mu=0$）化为
\[
R=\bar\partial\partial U.
\]

至此，我们引用 Yau 关于 $X$ 上具有指定 Ricci 曲率的 Kähler 度量存在性的定理。该定理断言：若 $c_1(X)=0$，则对 $X$ 上任一固定函数 $U$，方程 $R=\bar\partial\partial U$ 在给定 $X$ 上复结构与 Kähler 类的条件下有唯一解。现在，设 $\bar\omega_0$ 是 $X$ 上一个固定的 Kähler 类，我们试着对很大的 $p$、在 Kähler 类为 $p\bar\omega_0$ 的条件下解 $R=\bar\partial\partial U$。在首阶近似（$p=\infty$）下，我们干脆取 $R=0$ 的解。这是自洽的，因为当 $s\ge2$ 时，$U$ 的 $s$ 圈贡献在方程中给出相对阶为 $p^{-(s-1)}$ 的项。确定了首阶度量之后，就得到了 $U$ 的近似确定，然后尝试用逐次近似法解方程 $R=\bar\partial\partial U$。由 Yau 定理，满足该方程的 Kähler 度量在 $1/p$ 的所有有限阶上都存在且唯一。收敛性是另一回事；人们相信，若 $p$ 足够大——或者说，若 Kähler 类足够丰沛（ample）且远离 Kähler 锥（Kähler cone）的边界——这个过程会收敛，给出 $\beta$ 函数的一个零点。对今天的模型而言，这些条件只不过意味着 $-\operatorname{Im}(t)$ 应当足够大。

对很大的 $-\operatorname{Im}(t)$，理论的一个好的描述就是我们出发时的、映到 $X$ 的映射的 $\sigma$ 模型。今天我们要探讨的，实际上就是 $-\operatorname{Im}(t)$ 不大时会发生什么。

今天我们将看到 $c_1(X)=0$ 的 $X$ 的例子：存在一族由带 puncture 的柱面参数化的理论，它在 $\operatorname{Im}(t)\to-\infty$ 时趋近 $X$ 上的 $\sigma$ 模型，而在 $\operatorname{Im}(t)\to+\infty$ 时趋近一个 Landau--Ginzburg 模型。要得到这样一族理论，柱面必须在一点处被打洞（punctured），在那里理论有某种极点。人们相信这个带 puncture 的柱面参数化了一族光滑的理论，它们在红外流向一族超共形场论。作为这一信念的一个检验，我们论证 $+\infty$ 处 Landau--Ginzburg 理论的中心荷与 $-\infty$ 处 $\sigma$ 模型的中心荷相等。我们将在后面的一讲中为今天所描述的图景提供更多证据。

\subsubsection{考虑 $\mathcal{N}=2$ 超对称理论的理由}
% 原文 §13.2.2：Reasons for considering $\mathcal{N}=2$ SUSY theories

研究 $\mathcal{N}=2$ 超对称理论的一个好处是：它们比超对称更少（或没有超对称）的理论要刚性（rigid）得多。例如，若我们考虑映入紧 Riemann 流形 $X$ 的二维 $\sigma$ 模型，

% [Tutor Note: 「不变量理论表明任何闭 $(1,1)$-形式都是 $\bar\partial\partial U$」这一步作者一句话带过，实为整节的技术核心。可在正规坐标下把可用张量清点一遍：可用的只有 $g_{i\bar j}$、Riemann 张量 $R_{i\bar j k\bar l}$ 及其协变导数；次数为 $2$ 的闭不变量只有 Ricci 形式 $R_{i\bar j}$，而所有更高次数的闭不变量都可写成某个局部标量函数的 $\partial\bar\partial$——这正是「局部且普适」这个词的含义：构造不得依赖于流形上的额外结构（例如从固定环境空间拉回某个类），否则可以有非恰当的闭形式。附带提醒：原文印作 $R=\bar\partial\partial U$；由于 $\partial\bar\partial=-\bar\partial\partial$，与常见的 $\rho=\partial\bar\partial\log\det g$ 约定只差 $U$ 的一个符号。另需注意 Yau 定理给出的是「给定 Kähler 类」下的唯一性，故作者必须先固定 $\bar\omega_0$ 再谈 $p\bar\omega_0$；漏掉这个限定，方程 $R=\bar\partial\partial U$ 的解会在整个 Kähler 锥上滑动。]

% ===== 原文 p.176 =====
那么在无超对称时，我们可以向 Lagrange 量中加入形如
\[
\int d^2y\,\varphi^*(h)
\]
的项，其中 $h:X\to\mathbb{R}$ 是势函数。（这里 $\varphi:\Sigma\to X$ 是 Lagrange 量中的基本场。）在 $\mathcal{N}=1$ 超对称 $\sigma$ 模型中，这样的项是不可能的，但依然可以有
\[
\int d^2y\,d^2\theta\,\Phi^*(h)
\]
这样的项，其中 $h$ 仍是 $X$ 上的实值函数（称为超势），$\Phi$ 是 Lagrange 量中的基本超场。两种情形下这都是一个相关微扰。当我们要求 $\mathcal{N}=2$ 超对称时，Riemann 流形必须是 Kähler 的，而超势必须具有 $h=\operatorname{Re}(f)$ 的形式，其中 $f$ 是全纯函数。由于我们假定 Riemann 流形 $X$ 是紧的，这意味着 $f$ 为常数，从而在完成 $\theta$ 积分之后这一项为零。简单的量纲分析表明，在经典层面，这是 $2|4$ 超空间上唯一可能的相关微扰。这就是为什么在二维、$\mathcal{N}=2$ 超对称下，势为零的 $\sigma$ 模型是自然的；而在 $\mathcal{N}<2$ 时，它们既不自然、也更难研究。

\subsubsection{共形二维 $\sigma$ 模型的相关微扰与边际微扰}
% 原文 §13.2.3：Relevant and marginal perturbations of the conformal two-dimensional sigma model

让我们研究这些共形 $\sigma$ 模型的边际微扰（marginal perturbation）。$\sigma$ 模型的数据之一，就是 $X$ 上的复结构。改变复结构是该理论的一个边际微扰。让我们在局部坐标中考察这一微扰。$X$ 上的度量张量是形如 $h_{i\bar j}dz^i\otimes d\bar z^j$ 的 $(1,1)$-张量。复结构的改变对应于把度量作如下形式的微扰：
\[
\delta h_{\bar i\bar j}+\delta h_{ij},
\]
即 $(2,0)+(0,2)$ 型。当然，这个形式必须是实的，新的形式才能是度量张量；而 $\delta h_{i\bar j}$ 必须是 $\bar\partial$-闭的，新的近复结构（almost complex structure）才是可积的。让我们计算与 Kähler 度量的这一微扰相应的 Lagrange 量的改变。我们用 Riemann 面上的玻色坐标 $y^\alpha$（$\alpha=1,2$）与费米坐标 $\theta^\pm,\bar\theta^\pm$，在局部坐标中书写。超场记为 $X=X(y,\theta,\bar\theta)$，手征超场记为 $\Phi(y,\theta)$，反手征超场记为 $\bar\Phi(y,\bar\theta)$。Lagrange 量的改变由
\[
\Delta\mathcal{L}=\left(\int d^2y\,d^2\bar\theta\,\delta h_{i\bar j}\bar D_+X^{\bar i}\bar D_-X^{\bar j}\right)+\text{c.c.}
\]
给出。验证 $\Delta\mathcal{L}$ 中 $\bar\theta$ 积分的被积对象是反手征的，是一个很好的练习。（这里用到 $\delta h_{i\bar j}$ 是 $\bar\partial$-闭的这一事实。）这是两类可能的边际微扰之一。

另一类边际微扰，是固定复结构 $X$、而改变 Kähler 度量的复化 Kähler 类。这一形变的一部分，是用一个 Hermite 对称的闭二形式作通常的变分 $\delta h^{\mathrm{herm}}_{i\bar j}dz^{\bar i}\wedge dz^j$。这个通常的度量微扰是某个复微扰的实部，该复微扰的虚部称为「$B$ 场」（$B$-field）微扰（我们将会看到，它类似于规范理论中出现的 $\theta$ 角）。因此，我们

% [Tutor Note: 此处原文印作 $\delta h_{i\bar j}+\delta h_{ij}$，但按同句「$(2,0)+(0,2)$ 型」以及「此形式必须为实」的要求，正确的一对应当是 $\delta h_{ij}$（$(2,0)$）与 $\delta h_{\bar i\bar j}$（$(0,2)$）；译文中按几何意义写成 $\delta h_{\bar i\bar j}+\delta h_{ij}$，而把 $\delta h_{i\bar j}$ 留给 $(1,1)$ 型的 Kähler 类形变（即 p.177 的 $\delta h=\delta h_{i\bar j}d\bar z^i\wedge dz^j$）。这一区分并非小事：$(2,0)$ 部分改的是复结构（模空间上的复结构形变），$(1,1)$ 部分改的是 Kähler 类，二者是 $\sigma$ 模型模空间的两个不同因子，混淆指标会把两类边际微扰混为一谈。另请留意指标配对：下一段的 $\Delta\mathcal{L}$ 中 $\delta h$ 的第一个下标在原文里不带杠，而 $X$ 的两个上标都带杠（$X^{\bar i}$、$X^{\bar j}$）；按「同带杠性缩并」的规矩，这里的系数应当是 $\delta h_{\bar i\bar j}$（即同一形变的 $(0,2)$ 分量，其复共轭给出 $(2,0)$ 分量），故这一处横杠的遗漏同样只是排印疏漏。把两处对照起来看，原文对 $\delta h$ 的指标用法可归纳为：$(1,1)$ 型（Hermite）写 $\delta h_{i\bar j}$，复结构形变则须写成 $\delta h_{ij}$ 与 $\delta h_{\bar i\bar j}$ 这一对。]

% ===== 原文 p.177 =====
考虑微扰 $\delta h=\delta h_{i\bar j}d\bar z^i\wedge dz^j$，其中 $\delta h$ 是闭的复二形式。这导致 Lagrange 量的如下形式的微扰
\[
\Delta\mathcal{L}=\left(\int d^2y\,d\bar\theta^+d\theta^-(\delta h_{i\bar j}\bar D_+X^{\bar i}D_-X^j)\right)+\text{c.c.}
\]

除了上面描述的理论的边际形变之外，只可能有
\[
\int d^2y\,d^4\theta\,(\text{something})
\]
形式的形变。对它们而言，$\theta$ 积分的权为 $2$，因此若「something」是一个函数，则这类项在经典层面是边际的（维数为二）。（注意这个计数需要四个 $\theta$，因而用到了我们有 $\mathcal{N}=2$ 超对称这一事实。）但在量子层面，它们具有反常维数（anomalous dimension，正比于 Laplacian 的非零本征值，因而严格为正），从而是不相关的。这类形变可以在不改变复结构与 Kähler 类的情况下影响 Kähler 度量。

这些计算表明，对 $c_1(X)=0$ 的 $\sigma$ 模型，理论没有相关微扰，而边际微扰则由复结构与复化 Kähler 度量的形变给出。

让我们从这样一个 $\sigma$ 模型出发，其 Kähler 度量是光滑的、曲率半径很大，并且接近于（但不等于）真正给出零 $\beta$ 函数的那个度量。重整化群流会把我们带到给出共形不变性的度量，因为所有其它形变都是不相关的。假如存在相关微扰，它们在红外极限中会占主导，而那个极限一般来说就无法控制了。

\subsubsection{中心荷}
% 原文 §13.2.4：Central Charge

让我们回顾 [II-CFT, §2.4] 中关于二维共形与超共形场论中心荷的一些材料。在没有超对称时，我们可以把共形场论有效地分成两大类——中心荷 $c\le1$ 的，与中心荷 $c>1$ 的。$c<1$ 的理论称为极小模型（minimal models），只对离散的一组 $c$ 值存在，并可以用 Virasoro 代数的表示论以代数方式显式描述。$c=1$ 是 $c$ 的连续谱段的起点。$c=1$ 的理论也可以显式描述，用的是自由玻色子。当 $c\le1$ 时，左行与右行模式的 Virasoro 代数在量子 Hilbert 空间上几乎不可约地作用，这就是为什么这些理论被了解得相当清楚。对 $c>1$，Virasoro 作用远非不可约（最高权向量的数目随能量升高而指数增长），共形场论的味道就完全不同了。

当我们转到二维 $\mathcal{N}=1$ 超共形代数时，中心荷同样落在某个离散集合或某个连续区间之中。中心荷连续区间的起点，由自由标量超场所实现的理论给出。由于这个场是一个自由标量玻色子与一个自由标量费米子之和，若把它看作通常共形代数的表示，它具有中心荷 $c=3/2$。人们通常定义 $\tilde c=2c/3$，使得 $\tilde c$ 的连续值从 $\tilde c=1$ 开始。

% [Tutor Note: 这类「$\int d^4\theta\,K$」形变在量子层面的反常维数正比于 $X$ 上 Laplacian 的本征值，这一点值得补一句：紧流形上 $-\Delta$ 的谱非负，零模恰为常数，而常数 $K$ 对应的形变正是 Kähler 类的移动（也就是作者刚讨论过的 $(1,1)$ 边际形变）。因此真正的分界不是「$\theta$ 积分给权 $2$ 就算边际」，而是「$K$ 的零模边际、非零模不相关」：Laplacian 的非零本征值给出严格正的反常维数，把非零模一律推向不相关。这也解释了为何 $c_1(X)=0$ 的 $\sigma$ 模型只有复结构与复化 Kähler 类两类边际方向，其模空间是局部积结构。]

% ===== 原文 p.178 =====
对二维 $\mathcal{N}=2$ 超共形代数，情况是类似的。这里连续部分的最低值由自由手征超场实现。由于这个场由一个复标量玻色子与一个复费米子组成，它的 $\mathcal{N}=0$ 中心荷为 $3$。于是我们令 $\hat c=c/3$，使得连续部分的最低值再次是 $\hat c=1$。$\hat c<1$ 的 $\mathcal{N}=2$ 理论的离散集合就是「$\mathcal{N}=2$ 极小模型」，我们稍后会详细考虑它们。它们的中心荷取
\[
\hat c=1-\frac{2}{n},\qquad n=3,4,5,\dots
\]
的形式，并且可以用代数方式描述，也可以以多种诱人的方式加以研究。

靶空间为 $k$ 维复流形 $X$、且 $c_1(X)=0$ 的共形 $\mathcal{N}=2$ 超对称 $\sigma$ 模型，其中心荷为 $\hat c=k$。证明方式是利用 $\hat c$ 在一族共形场论中为常数这一事实：把度量不断展开，直到 $X$ 在局部区域可以被带平直度量的 $\mathbb{C}^k$ 以任意精度逼近。于是这个 $\sigma$ 模型与靶为平直 $\mathbb{C}^k$ 的 $\sigma$ 模型有相同的 $\hat c$。后者是含 $k$ 个手征超场的自由理论，因此 $\hat c=k$，如所断言。

% [Tutor Note: $\hat c=1-2/n$ 与标准 $\mathcal{N}=2$ 极小模型的公式可以对上，是一个很好的自查点：后者常写成 $c=\frac{3k}{k+2}$（$k=1,2,\dots$ 为 ADE 型级别），于是 $\hat c=c/3=\frac{k}{k+2}=\frac{n-2}{n}$，只需令 $n=k+2$ 即得 $n=3,4,5,\dots$；$n=3$ 给出 $\hat c=1/3$（即 $A_1$ 模型，等价于一个自由玻色子的轨形）。这也提醒：原文的 $n$ 是从 $3$ 起计数的，不要与超场个数 $n$ 混同——本讲后面 §13.4 的 $n$ 恰好是手征超场个数，两个 $n$ 没有任何关系。]

\subsection{$U(1)$ 理论}
% 原文 §13.3：The $U(1)$ Theories

现在让我们转向研究 $U(1)$ 规范理论，带有复值手征超场 $A_1,\dots,A_n$。这些超场在 $U(1)$ 下带荷，荷为 $q_1,\dots,q_n\in\mathbb{Z}$。回想 [I-Supersolutions] 第 9 章：我们在 $2|4$ 超空间上所用的联络，被约束为在所有的奇方向对上都平坦，唯一例外是 $\{\bar\theta^+,\theta^-\}$ 这一对（及其复共轭）。这样一个联络的基本不变量，是奇方向上曲率的非零部分：
\[
\Sigma=\{\bar{\mathcal{D}}_+,\mathcal{D}_-\}.
\]
$\Sigma$ 是一个扭曲手征超场（twisted chiral superfield），取值于 $2|4$ 超空间上伴随丛（adjoint bundle）的复化（complexification）之中。我们能写下的最小纯规范理论 Lagrange 量是
\[
\int d^2y\,d^4\theta\,\frac{1}{2e^2}\bar\Sigma\Sigma.
\]
再加上 $n$ 个手征超场的动能项后，最小 Lagrange 量成为
\begin{equation}
  \tag{13.1}
  \label{eq:13-1}
  \mathcal{L}_0=\int d^2y\,d^4\theta\left(\frac{1}{2e^2}\bar\Sigma\Sigma+\sum_{i=1}^{n}\bar A_iA_i\right).
\end{equation}
最终我们还会向这个 Lagrange 量加入一个超势项与一个扭曲超势项。

\subsubsection{$R$ 对称性}
% 原文 §13.3.1：$R$-symmetries

回想 [I-Supersolutions] 第 9 章：上述二维 $\mathcal{N}=2$ 超对称规范理论（一个 $2|2$ 超理论）来自四维 $\mathcal{N}=1$ 规范理论（一个 $4|4$ 超理论）的维数约化。因此，二维的 $\mathcal{N}=2$ 理论有一个由四维空间中剩余两个方向上的空间转动所诱导的 $R$ 对称性。但此外，四维中 $U(1)$ 规范理论的 $\mathcal{N}=1$ 超对称版本本身也有一个 $R$ 对称性。这个来自四维的 $R$ 对称性经维数约化，给出我们二维理论的第二个 $R$ 对称性。

% ===== 原文 p.179 =====
不管它们的来源如何，在 $2|4$ 理论中直接看出 $U(1)\times U(1)$ 这组 $R$ 对称性是什么，是容易的。我们有 $J_R$，即圆周作用
\[
\begin{gathered}
\theta^+\mapsto e^{i\alpha}\theta^+,\\
\bar\theta^+\mapsto e^{-i\alpha}\bar\theta^+,
\end{gathered}
\]
其中 $\theta^-$ 保持不动；我们还有 $J_L$，即作用在 $\theta^-$ 上、$\theta^+$ 保持不动的那一个圆周作用。来自四维规范理论 $R$ 对称性的那个 $R$ 对称性，是这两者的乘积（也就是说，它在两个 $\theta^\pm$ 上都按标准表示作用）。被省略的那两个方向中的空间转动所诱导的圆周作用，则是第一个作用乘以第二个作用的逆。

$R$ 对称性很重要，因为 $\mathcal{N}=2$ 超共形代数包含左、右 $R$ 对称性。因此每个 $\mathcal{N}=2$ 超共形理论都有左、右 $R$ 对称性。我们经典理论的其中一个 $R$ 对称性可能带有反常。若我们的理论要在红外流向一个超共形理论、且该理论的超共形代数中包含这些 $R$ 对称性，则反常必须为零。（也可能流向一个超共形场论，其不动点理论的 $R$ 对称性不是原 Lagrange 量的对称性；但在那种情形下，Lagrange 量描述往往不那么有用。）

让我们把先前关于 $\sigma$ 模型的反常正比于第一 Chern 类 $c_1(X)$ 的评述说得更明确些。要点在于：当我们把 Lagrange 量用坐标写出来时，得到的是形如
\[
\mathcal{L}=\int d^2y\left(g_{i\bar j}\partial X^i\bar\partial X^{\bar j}+\bar\psi^{(0,1)}_-\not{\partial}_+\bar\psi^{(1,0)}_-+\psi^{(0,1)}_+\not{\partial}_-\psi^{(1,0)}_+\right)
\]
的表达式（外加一个不影响当前讨论的四费米子项）。

% [Tutor Note: 第一组费米子动能项在原文中印作 $\bar\psi^{(0,1)}_-\not{\partial}_+\bar\psi^{(1,0)}_-$，两个因子都带横杠；按动能项必须是「$\bar\psi$ 与 $\psi$ 配对」的规矩，第二个横杠几乎可以肯定是排印时多打的一个（对照第二组 $\psi^{(0,1)}_+\not{\partial}_-\psi^{(1,0)}_+$ 就完全不带横杠）。译文按原文照录公式本体，但读者应按 $\bar\psi^{(0,1)}_-\not{\partial}_+\psi^{(1,0)}_-$ 理解，否则这一项的费米子数不为零。附带说明记号：括号中的 $(0,1)$ 与 $(1,0)$ 是取值丛的 Dolbeault 型，下标 $\pm$ 是世界面手征性；$\not{\partial}_\pm$ 分别作用在 $(1,0)$-值与 $(0,1)$-值旋量上，因此「$\psi^{(0,1)}$ 配 $\not{\partial}_+$、$\psi^{(1,0)}$ 配 $\not{\partial}_-$」是固定的配对方式，不能随意调换，否则下面 $\operatorname{index}(\not{\partial}_\pm)$ 的符号就无从判断。]

两个基本 $R$ 对称性之一是 $J_L$，它转动 $\psi^{(0,1)}_-$ 与 $\psi^{(1,0)}_-$；另一个是 $J_R$，它以同样的方式转动 $\psi^{(0,1)}_+$ 与 $\psi^{(1,0)}_+$。于是，在由 $a_RJ_R+a_LJ_L$ 生成的手征转动下，路径积分中的测度 $\mathcal{D}\psi_+\mathcal{D}\psi_-$ 的改变量正比于
\[
a_R\operatorname{index}(\not{\partial}_+)+a_L\operatorname{index}(\not{\partial}_-).
\]
算符 $\not{\partial}_\pm$ 是正、负手征性的 Dirac 算符，作用在 $X$ 的全纯切丛 $T^{(1,0)}(X)$ 拉回到 Riemann 面上的那个丛上。由指标定理，这两个算符的指标都由 $c_1(X)$ 拉回、并在 Riemann 面的基本类上求值给出。于是我们看到，来自四维 $R$ 对称性的那个 $R$ 对称性没有反常，因为 $\operatorname{index}(\not{\partial}_+)-\operatorname{index}(\not{\partial}_-)=0$；但来自空间转动的那个 $R$ 对称性却可以有反常：$\operatorname{index}(\not{\partial}_+)+\operatorname{index}(\not{\partial}_-)=2\operatorname{index}(\not{\partial}_+)$。这一反常对所有映入 $X$ 的映射同伦类、对所有曲面 $\Sigma$ 都为零，当且仅当 $c_1(X)=0$。

我们刚刚给出了 $\sigma$ 模型反常的计算。让我们在带手征超场 $A_i$（$1\le i\le n$）的 $U(1)$ 规范理论中作类似的计算。超场 $A_i$ 是与 $2|4$ 超空间上联络相联系的那个线丛的 $q_i$ 次幂的截面。它可以按通常的场（玻色子与费米子）展开。$R$ 对称性在费米子上作为手征对称性作用。因此与上面同类型的「反常—指标」计算表明，$J_R$ 的反常由
\[
\sum_i q_i\int d^2y\,\frac{F_A}{2\pi}
\]
给出。

% [Tutor Note: 这里两句话的归属与它们前面的公式对不上，属于原文（学生笔记）的一处笔误，宜按公式为准重读。按公式，反常系数是 $a_R\operatorname{index}(\not{\partial}_+)+a_L\operatorname{index}(\not{\partial}_-)$：于是 $a_R=a_L$（即 p.179 上所说「在两个 $\theta^\pm$ 上按标准表示作用」的乘积 $J_RJ_L$，来自四维 $R$ 对称性）对应和式 $\operatorname{index}(\not{\partial}_+)+\operatorname{index}(\not{\partial}_-)=2\int_\Sigma\varphi^*c_1(X)$，一般不为零；而 $a_R=-a_L$（即 $J_RJ_L^{-1}$，来自空间转动）对应差式，因为原文自己指出两个指标都等于 $\int_\Sigma\varphi^*c_1(X)$，故差式恒为零。也就是说，「无反常」与「有反常」两句的归属应互换。结论（除非 $c_1(X)=0$，总有一个 $R$ 对称性反常）不受影响，这也正是二维 $(2,2)$ 理论中「有反常的那一个 $R$ 对称性恰好给出 $\beta$ 函数零点处 $c_1=0$ 的条件」的标准表述。另请注意 $\not{\partial}$ 的写法：$J_L$ 转动的是下标带「$-$」的费米子，而与之配对的 Dirac 算符是 $\not{\partial}_-$，与 $J_R$ 配对的才是 $\not{\partial}_+$，读时勿按符号顺序误配。]

% ===== 原文 p.180 =====
其中 $F_A$ 是 Riemann 面上通常联络的曲率。（$J_L$ 的反常是这个数的相反数。）因此，要得到超共形场论，我们需要
\[
\sum_i q_i=0.
\]

\subsection{一个例子}
% 原文 §13.4：One Example

为使事情尽可能简单具体，让我们聚焦于 $U(1)$ 规范理论的一个特例：$n+1$ 个手征超场
\[
\begin{gathered}
A_1,\dots,A_n\quad\text{的权 }q=1,\\
A_{n+1}=P\quad\text{的权 }q=-n.
\end{gathered}
\]
这样选择权，保证了 $\sum_iq_i=0$，从而我们的 $R$ 对称性无反常。现在是时候用超势与扭曲超势来补充式 13.1 的最小 Lagrange 量 $\mathcal{L}_0$ 了。也就是说，我们考虑形如
\begin{equation}
  \tag{13.2}
  \label{eq:13-2}
  \mathcal{L}=\mathcal{L}_0+\left(\int d^2y\,d\theta^+d\theta^-\,W(A,P)+\text{c.c.}\right)+\left(\int d^2y\,d\bar\theta^+d\theta^-\,\widetilde W(\Sigma)+\text{c.c.}\right)
\end{equation}
的 Lagrange 量。这里，$W$ 称为超势；它必须是 $\mathbb{C}^{n+1}$ 上的全纯函数。函数 $\widetilde W(\Sigma)$ 称为扭曲手征超势项。

我们希望选取 $W$ 与 $\widetilde W$ 以保持 $R$ 对称性。先考虑扭曲手征超势项 $\widetilde W$。这里，利用我们的理论是由 $4|4$ 空间上的规范理论经维数约化而得到这一事实，是有帮助的。从 $4|4$ 空间上的联络 $A=\sum_{\mu=0}^{3}A_\mu dx^\mu$ 出发，我们构造
\[
\sum_{\alpha=0}^{1}dy^\alpha+\sigma(dx^2+idx^3)+\text{c.c.}.
\]
这告诉我们 $\Sigma=\sigma+\dots$ 在 $R$ 对称性下如何变换：它的变换方式与 $dx^2+idx^3$ 相反。要使式 13.2 中来自扭曲超势的额外项不变，它必须在 $R$ 对称性下具有与 $\Sigma$ 相同的变换规律。也就是说，$\widetilde W(\Sigma)$ 必须是 $\Sigma$ 的线性函数，即对某个 $t\in\mathbb{C}$ 有 $\widetilde W(\Sigma)=t\Sigma$。记
\[
t=\frac{\theta}{2\pi}-ir,
\]
我们便有
\[
\left(\int d^2y\,d\bar\theta^+d\theta^-\,t\Sigma\right)+\text{c.c.}=\theta\int\frac{F}{2\pi}+r\int D.
\]
由此可见，加入这样一个线性的扭曲超势，会产生一个 Fayet--Iliopoulos 项 $r\int D$，以及一个 $\theta$ 角项，后者是 $\theta$ 乘以某个积分拓扑量。Fayet--Iliopoulos 项对应于给 $\mathbb{C}^{n+1}$ 上 $U(1)$ 作用的矩映射（moment map）加上一个常数。

现在我们转向超势 $W(A,P)$。为使
\[
\int d^2y\,d\theta^+d\theta^-\,W(A,P)
\]
这一项不变，我们需要 $W(A,P)$ 在由四维 $R$ 对称性诱导的 $R$ 对称性下以权 $2$ 变换，并且在由四维空间转动诱导的 $R$ 对称性下不变。为了安排这一点，

% [Tutor Note: 此处 $t$ 的写法与 p.173 的 $t=-ir+\theta/2\pi$ 完全一致，但符号最容易出错的是下一行：$t\Sigma$ 的 $\theta$ 积分给出 $\theta\int F/2\pi+r\int D$，即 Fayet--Iliopoulos 系数取 $+r$ 而不是 $-r$，因为 $\theta$ 积分在 $\Sigma$ 上取出的正是 $D$ 项、而 $t$ 的虚部是 $-ir$。若把 $t$ 写成文献中常见的 $\theta/2\pi+ir$，这一行的两个符号会同时翻转，而 $\theta\to\theta+2\pi$、$r\to r$ 的周期性仍成立，故单看周期性无法判定约定，必须回到 $t=-ir+\theta/2\pi$ 这个定义。另请注意：$\Sigma$ 由 $\{\bar{\mathcal D}_+,\mathcal D_-\}$ 定义，故它在 $J_R$ 下与 $\bar\theta^+$ 同权、在 $J_L$ 下与 $\theta^-$ 同权；只有线性 $\widetilde W=t\Sigma$ 能同时满足「权为零」与「$R$ 对称性不变」，这正是本讲反复强调的 $\mathcal{N}=2$ 刚性——扭曲超势没有高次项可加（对比之下超势 $W$ 只要在相应权下总次数为 $2$ 就可以是高次多项式）。最后，「$W$ 以权 $2$ 变换」这一条件结合 $q=(1,\dots,1,-n)$ 就是下一讲反复使用的「权条件」（weight condition）的具体形式：它要求 $W$ 为拟齐次（quasi-homogeneous）多项式、在 $R$ 对称性下总权恰为 $2$；而 $\sum_iq_i=0$ 保证了这一 $R$ 对称性无反常。请注意权条件与规范不变性是两个不同的要求：按这些荷，$W=P\,f_n(A)$（$f_n$ 为 $A$ 的 $n$ 次齐次多项式）的规范荷为 $-n+n=0$，故规范不变；它的 $R$ 权为 $2$ 则是把 $R$ 荷取成规范荷的适当平移（$\sum_iq_i=0$ 使这种平移成为可能）之后得到的。阅读时不要把这两个条件混为一谈。]
